\section{An \texorpdfstring{\(\Ein\)}{Ein}-linear realization of
characteristic-zero algebraic matroids}
\label{app:Ein-matroids}

This appendix is the only place where the paper uses the linearization
theorem of Rosen--Sidman--Theran.  It is kept separate so that the moving
trace and Ramanujan results do not appear to depend on it.

\begin{theorem}[\(\Ein\)-realization of algebraic matroids]
\label{thm:matroid-Ein-realization}
Let \(M\) be a finite characteristic-zero algebraic matroid of rank \(r\)
on a ground set \(E\).  There are a matrix
\(B=(b_{je})\in\Qbar^{r\times E}\) representing \(M\), distinct
\(\beta_1,\ldots,\beta_r\in\Qbar^\times\), and numbers
\begin{equation}\label{eq:matroid-Ein-realization}
 \xi_j=\Ein(\beta_j),\qquad
 W_e=\sum_{j=1}^r b_{je}\xi_j,
\end{equation}
such that, for every \(S\subseteq E\),
\begin{equation}\label{eq:matroid-rank-realization}
 \boxed{\displaystyle
 \trdeg_KK(W_e:e\in S)=r_M(S).}
\end{equation}
Moreover the complete relation ideal is
\begin{equation}\label{eq:matroid-linear-ideal}
 \boxed{\displaystyle
 \ker\!\left(K[X_e:e\in E]\longrightarrow K[W_e:e\in E]\right)
 =\left\langle\sum_{e\in E}c_eX_e:
 c\in\Qbar^E,\ Bc=0\right\rangle.}
\end{equation}
\end{theorem}

\begin{proof}
Extend a characteristic-zero field of algebraic realization to its
algebraic closure; the algebraic-dependence matroid is unchanged.  By
\cite[Thm.~1.10]{RosenSidmanTheran2025}, the resulting matroid is linearly
representable over that algebraically closed field.  Linear
representability of a fixed finite
matroid is described by the vanishing of its nonbasis minors and the
nonvanishing of its basis minors.  Introducing one auxiliary inverse for
the product of the basis minors converts these conditions into polynomial
equations over \(\Q\).  Since they have a solution in an algebraically
closed characteristic-zero field, the Nullstellensatz gives a solution in
\(\Qbar\).  Thus a representing matrix \(B\) may be chosen over \(\Qbar\).

Choose distinct nonzero algebraic \(\beta_j\).  By
\cref{thm:input-independence}, the \(\xi_j\) are algebraically independent
over \(K\).  If the columns of \(B\) indexed by \(S\) have rank \(q\),
choose \(q\) independent columns.  The corresponding \(W_e\) are
algebraically independent by \cref{lem:linear-images}, giving the lower
bound \(q\); every remaining column is a \(\Qbar\)-linear combination of
those columns, giving the upper bound.  This proves
\eqref{eq:matroid-rank-realization}.  Finally the symmetric algebra of the
column span is the polynomial ring on the \(X_e\) modulo the linear forms
coming from \(\ker B\), which proves \eqref{eq:matroid-linear-ideal}.
\end{proof}

\begin{remark}[Exact scope]
The construction realizes the rank function of \(M\) and the linearized
relation ideal \eqref{eq:matroid-linear-ideal}.  It does not reproduce the
nonlinear circuit polynomials of a chosen algebraic presentation of
\(M\), and it is not a claim that such a presentation becomes linear as an
algebraic variety.
\end{remark}

\section{Creative-telescoping certificates for the Euler recurrence}
\label{app:telescopers}

For \(a=n+3\), \(b=n+4\), define meromorphically
\[
 T_{n,k}=b\frac{(a!b!)^2}
 {\Gamma(k+1)^3\Gamma(a-k+1)\Gamma(b-k+1)},
 \qquad
 \mathcal L_{n,k}=3H_k-H_{a-k}-H_{b-k},
\]
where \(H_z=\psi(z+1)+\gamma\).  The regularized endpoint values are
\[
 T_{n,b}\mathcal L_{n,b}=a!,
 \qquad
 T_{n,k}\mathcal L_{n,k}=0\quad(k>b).
\]
Indeed, at \(k=b\) the simple zero of
\(1/\Gamma(a-k+1)\) cancels the simple pole of
\(-H_{a-k}\), with limiting product \(a!\).  For \(k>b\), the two
reciprocal-gamma zeros dominate the at-most-simple digamma pole, so the
regularized product vanishes.
Thus \eqref{eq:qn-closed} and \eqref{eq:pn-closed} are respectively
\(\sum_{k\ge0}T_{n,k}\) and
\(\sum_{k\ge0}T_{n,k}\mathcal L_{n,k}\).

Let
\[
 c_0=-A_n,
 \qquad c_1=B_n,
 \qquad c_2=-C_n,
 \qquad c_3=D_n,
\]
and \(\Delta_ku(k)=u(k+1)-u(k)\).  The certificates are
\begin{align}
 \sum_{j=0}^3c_jT_{n-j,k}
 &=\Delta_k\{T_{n,k}h_{n,k}\},
 \label{eq:q-telescoper}\\
 \sum_{j=0}^3c_jT_{n-j,k}\mathcal L_{n-j,k}
 &=\Delta_k\{T_{n,k}(h_{n,k}\mathcal L_{n,k}+g_{n,k})\}.
 \label{eq:p-telescoper}
\end{align}
Here
\[
 h_{n,k}=-\frac{k^3}{(n+3)^3(n+4)^3}\mathcal H(n,k),
 \qquad
 g_{n,k}=\frac{k^2}{(n+3)^3(n+4)^3}\mathcal G(n,k),
\]
where
\begin{align*}
\mathcal H(n,k)={}&
 (8n^3+75n^2+231n+232)k^3\\
&-(48n^4+562n^3+2436n^2+4626n+3248)k^2\\
&+(96n^5+1388n^4+7924n^3+22278n^2+30772n+16668)k\\
&-(64n^6+1120n^5+8071n^4+30590n^3+64148n^2+70347n+31388),
\end{align*}
and
\begin{align*}
\mathcal G(n,k)={}&
 (48n^3+450n^2+1386n+1392)k^3\\
&-(240n^4+2810n^3+12180n^2+23130n+16240)k^2\\
&+(384n^5+5552n^4+31696n^3+89112n^2+123088n+66672)k\\
&-(192n^6+3360n^5+24213n^4+91770n^3+192444n^2+211041n+94164).
\end{align*}

To verify \eqref{eq:q-telescoper}, divide by \(T_{n,k}\), use
\[
 \frac{T_{n,k+1}}{T_{n,k}}
 =\frac{(n+3-k)(n+4-k)}{(k+1)^3},
\]
and, for \(0\le j\le3\),
\[
 \frac{T_{n-j,k}}{T_{n,k}}
 =\frac{n-j+4}{n+4}
 \frac{(n+3-k)^{\underline j}(n+4-k)^{\underline j}}
 {\{(n+3)^{\underline j}(n+4)^{\underline j}\}^{2}},
\]
where \(x^{\underline j}=x(x-1)\cdots(x-j+1)\), and clear denominators.
For \eqref{eq:p-telescoper}, also use
\[
 \mathcal L_{n,k+1}-\mathcal L_{n,k}
 =\frac3{k+1}+\frac1{n+3-k}+\frac1{n+4-k}.
\]
The remaining shifts are reduced by
\[
 \mathcal L_{n-j,k}-\mathcal L_{n,k}
 =\sum_{r=0}^{j-1}
 \left(\frac1{n+3-k-r}+\frac1{n+4-k-r}\right).
\]
Both become polynomial identities in \(n,k\).  The lower boundary vanishes
because of the factors \(k^3,k^2\), and the reciprocal gamma factors give
the upper boundary.  Summing proves \eqref{eq:challenge-recurrence} for both
closed forms for every \(n\ge0\).  The accompanying source package contains
an exact symbolic verifier for these identities and for
\eqref{eq:cmf-det}--\eqref{eq:cmf-flatness}.

\section{Finite certificate for the integer initial lattice}
\label{app:integer-lattice-certificate}

The official dual transition and cyclic gauge may be taken as
\begin{equation}\label{eq:Rn-certificate}
 \mathsf R(n)=
 \begin{pmatrix}
 6n^2+21n+18&-8n-11&3\\
 3n^3+12n^2+16n+7&-3n^2-7n-4&n+1\\
 (n+1)^4&0&0
 \end{pmatrix}
\end{equation}
and
\begin{equation}\label{eq:Un-certificate}
\mathsf U(n)=
\begin{pmatrix}
1&6n^2+21n+18&15n^4+147n^3+520n^2+792n+440\\
0&3n^3+12n^2+16n+7&10n^5+96n^4+353n^3+625n^2+537n+179\\
0&(n+1)^4&6n^6+57n^5+213n^4+402n^3+408n^2+213n+45
\end{pmatrix}.
\end{equation}
Let the row-companion transfer of \eqref{eq:challenge-recurrence} be
\begin{equation}\label{eq:row-companion-certificate}
 \mathsf C_n^{\mathrm{comp}}=
 \begin{pmatrix}
 0&0&D_n/A_n\\
 1&0&-C_n/A_n\\
 0&1&B_n/A_n
 \end{pmatrix}.
\end{equation}
Direct multiplication gives the exact coboundary
\begin{equation}\label{eq:RU-companion-coboundary}
 \mathsf R(n)\mathsf U(n+1)
 =\mathsf U(n)\mathsf C_n^{\mathrm{comp}}.
\end{equation}
For a raw row \(c\), the integer recurrence state at level \(N\) is
\begin{equation}\label{eq:raw-state-certificate}
 c\,\mathsf R(0)\cdots\mathsf R(N-1)\mathsf U(N).
\end{equation}
Indeed \eqref{eq:RU-companion-coboundary} proves by induction that
\eqref{eq:raw-state-certificate} equals
\((u_{N-3},u_{N-2},u_{N-1})\) when \(c\mathsf U(0)\) is the initial row;
the empty product is used at \(N=0\).
Let
\[
 d_1=(8,-22,10),\qquad d_2=(0,2,-6),\qquad d_3=(0,-1,7).
\]
Direct multiplication gives
\begin{align}
 \frac{d_1}{8}\mathsf U(0)&=(1,0,4),&
 \frac{d_2}{8}\mathsf U(0)&=(0,1,11),&
 \frac{d_3}{8}\mathsf U(0)&=(0,0,17),
 \label{eq:lattice-basis-certificate}
\end{align}
and the finite congruence certificate
\begin{align}
 \mathsf R(0)\cdots\mathsf R(7)&\equiv0\pmod8,
 \label{eq:R8-zero}\\
 d_i\mathsf R(0)\cdots\mathsf R(N-1)\mathsf U(N)&\equiv0\pmod8
 \quad(1\le i\le3,\ 0\le N<8).
 \label{eq:finite-mod8-certificate}
\end{align}
For \(N\ge8\), \eqref{eq:R8-zero} supplies the factor eight; for
\(N<8\), \eqref{eq:finite-mod8-certificate} does.  Hence the three rows in
\eqref{eq:lattice-basis-certificate} generate integral solutions at all
levels, completing the sufficiency argument in
\cref{thm:ram-integer-lattice}.  The source verifier checks every displayed
matrix identity over \(\Z\).

\section{A numerical ledger}
\label{app:numerics}

The following values provide reproducibility checks for the indexing in
\eqref{eq:dN-def}.  They are not used in any proof.

\begin{table}[ht]
\centering
\caption{Positive increments and positive zeros of \(\Delta_2\).}
\begin{tabular}{@{}rcc@{}}
\toprule
\(N\)&\(d_N\)&\(r_N=\sqrt{2/d_N}\)\\
\midrule
0&1.474569389580200&1.1646150303\\
1&\(5.633945512124626\times10^{-1}\)&1.8841205074\\
2&\(9.586804440833174\times10^{-2}\)&4.5674948216\\
3&\(1.339700176289053\times10^{-3}\)&38.637693279\\
4&\(2.954643417466919\times10^{-7}\)&2601.7314103\\
5&\(2.050814843583590\times10^{-14}\)&\(9.8753335\times10^6\)\\
\bottomrule
\end{tabular}
\end{table}

The indexing is tied to the convention \(Y_N\) for \(N\ge1\):
\[
 d_0=\tau_1,\qquad d_N=\tau_{N+1}-\tau_N\quad(N\ge1).
\]
Starting a level sequence at \(N=0\) would give a different, though equally
valid, spectral decomposition.
